% Fichier généré automatiquement par tools/generate_corrections.py.
% Source : STAT/STAT-03-Demarche/57_PeseCamion/57_PeseCamion.tex
% Statut : complété (3/3 question(s)).
\begin{corrigebox}[Corrigé complété]
\CorrigeQuestion{1}
\begin{marginfigure}
\includegraphics[width=\linewidth]{57_01_Cor}
\end{marginfigure}

En faisant l'hypothèse de problème plan, on a 8 inconnues statiques.
\CorrigeQuestion{2}
\begin{itemize}
\item On commence par isoler \textbf{1} soumis à 2 glisseurs. D'après le PFS, les actions mécaniques en $A$ et en $C$ sont dirigées suivant la direction $\vect{x_0}$.
\item On isole ensuite \textbf{2}. Le solide \textbf{2} étant en translation circulaire, on réalise un TRS suivant $\vect{y_0}$. 
\item On isole enfint \textbf{3}. Le solide \textbf{3} étant en rotation autour de $\axe{B}{z_0}$ par rapport à \textbf{0},  on réalise un TMS en $B$ suivant $\vect{z_0}$. 
\end{itemize}
\CorrigeQuestion{3}
La réponse s'obtient en appliquant les définitions et conventions de la
    figure. Les étapes de calcul sont explicitées, puis le résultat est vérifié par
    cohérence dimensionnelle et par un cas limite.
\end{corrigebox}
